\midsectiontitle{Terrenos em Combate}

No mundo de Overgrown terrenos são parte importante do combate:

Movimento do jogador no sistema:
5 + 1 Grace até o limite de 15 metros

Tipos de terreno:
Abrasado: Terreno em chamas, alvos sobre terreno abrasado sofrem 1D8 de dano no inicio da rodada
\originlink{terr:alagado}
Alagado: Terreno alagado, todos os alvos dentro ficam vulnerareis a magias de eletricidade
\originlink{terr:congelado}
Congelado: Terreno congelados afetam o movimento, toda rodada que o alvo andar é necessário fazer um teste de Grace com atletismo, DT 15, se falhar ficará \link{status:derrubado}{Derrubado}.
Comprometido: Terrenos comprometidos são caracterizados por dificultar andar, reduzem o movimento do jogador pela metade, até sair.
Profanado: Terreno repleto de energia negativa, está presente nele consome energia nula.
Sagrado: Terreno magico criado pelos deuses, ficar em locais sagrados regenera HP passivamente
Nulo: Terreno que emana IEP ou energia nula, ficar em locais nulo, regenera IEP passivamente.
Nevoa Espessa: Terreno com Baixa visibilidade, desvantagem em todos os teste de ataque e percepção, vantagem em testes de furtividade
Frágil: Terreno sobrevoando abismo ou vale, pode ser quebrado com ataques físicos para destruir.